%\# -*- coding: utf-8-unix -*-

\chapter{Siegel-Shidlovsky 定理}\label{ch:11}

\section{引言}\label{sec:ch11_1}

1929年，Siegel\index{Siegel, C. L.}\footnote{Abh. Preuss. Akad. Wiss. 1929, No. 1.} 获得了一种建立满足线性微分方程组的某类幂级数值的代数无关性的通用方法. Siegel\index{Siegel, C. L.} 将所讨论的幂级数称为 $E$-函数. 他的意思是形如
\[\sum_{n=0}^{\infty} a_n x^n / n!,\]
的级数，其中 $a_0, a_1, \ldots$ 为代数数\index{algebraic number}域中的元素，使得对于某正整数序列 $b_0, b_1, \ldots$ 以及任意的 $\epsilon > 0$，$b_n a_0, \ldots, b_n a_n$ 和 $b_n$ 均为代数整数\index{algebraic integer}，其大小 $\leqslant n^{\epsilon n}$，其中隐含的常数仅依赖于 $\epsilon$；此处的大小如\autoref{ch:4}中所述，表示共轭元绝对值的最大值. 显然，指数函数\index{exponential function}是一个 $E$-函数，同样，标准化的 Bessel 函数\index{Bessel function}
\[K_{\lambda}(x) = \Gamma(\lambda+1) (\frac{1}{2}x)^{-\lambda} J_{\lambda}(x) = \sum_{n=0}^{\infty} \frac{(-1)^n (\frac{1}{2}x)^{2n}}{n! (\lambda+1) \dots (\lambda+n)}\]
对于除负整数外的所有有理数 $\lambda$ 也是一个 $E$-函数. 更一般地，任何超几何级数\index{hypergeometric function}
\[\sum_{n=0}^{\infty} \frac{[\alpha_1, n] \dots [\alpha_l, n]}{[\beta_1, n] \dots [\beta_m, n]} x^{kn}\]
是一个 $E$-函数，其中 $k = m - l > 0$，
\[[\gamma, n] = \gamma(\gamma+1) \dots (\gamma+n-1),\]
且所有 $\alpha$ 和 $\beta$ 均为除负整数外的有理数. 后一个断言实际上源于以下观察：对于任何有理数 $\alpha = p/q$，整数 $q^n[\alpha, n]$ 整除 $n! \nu$，其中 $\nu$ 表示所有不超过
\[m = (|p| + |q|) n\]
的正整数的最小公倍数；并且由素数定理，对于某个绝对常数 $c$，我们有 $\nu < c^m$. 此外，极易验证 $E$-函数的和、积、导数与积分依然是 $E$-函数.

Siegel\index{Siegel, C. L.} 的工作仅涉及一阶和二阶微分方程，并且许多年来，设计一种将这些论证推广到高阶方程的方法一直是一个悬而未决的问题. 该问题由 Shidlovsky\index{Shidlovsky, A. B.}\footnote{I.A.N. 23 (1959), 35--66.} 于1954年解决，随后出现了许多显著的应用. 基本结果涉及满足齐次线性微分方程组
\[y'_{i} = \sum_{j=1}^{n} f_{ij}(x) y_{j} \quad (1 \leq i \leq n), \tag{1}\label{eq:ch11_1}\]
的 $E$-函数 $E_1(x), \dots, E_n(x)$，其中 $f_{ij}$ 是 $x$ 的有理函数，并且所有 $E$ 和 $f$ 的系数均假定为代数数\index{algebraic number}域 $K$ 中的元素. 我们于是有：

\begin{mythm}{}{ch11_1}
若 $E_1(x), \dots, E_n(x)$ 在 $K(x)$ 上代数无关，则对于任何不同于 $f_{ij}$ 极点的非零代数数\index{algebraic number} $\alpha$，$E_1(\alpha), \dots, E_n(\alpha)$ 代数无关.
\end{mythm}

该定理很容易推广，从而得到一个断言，其大意是在 $E_1(x), \ldots, E_n(x)$ 中的最大代数无关元个数与在 $E_1(\alpha), \ldots, E_n(\alpha)$ 中的相同\footnote{参见 Feldman\index{Feldman, N. I.} 和 Shidlovsky\index{Shidlovsky, A. B.} 的综述（文献目录）. }，并且将后一个结果推广到右侧存在附加有理函数的非齐次方程并无困难. 作为\autoref{thm:ch11_1}的直接应用，我们看到如果 $\lambda$ 为有理数，但不是负整数或半奇数，则对于每个非零代数数\index{algebraic number} $\alpha$，$K_{\lambda}(\alpha)$ 和 $K'_{\lambda}(\alpha)$ 代数无关；因为众所周知\footnote{参见 Siegel\index{Siegel, C. L.}（文献目录）. }，$K_{\lambda}(x)$ 和 $K'_{\lambda}(x)$ 在 $\mathbb{Q}(x)$ 上代数无关. 这进一步说明，例如，偏商为 1, 2, 3, \ldots{} 的连分数是超越的；因为 $J_0(\sqrt{-4x})$ [$= K_0(\sqrt{-4x})$] 满足微分方程 $xy'' + y' = y$，并且所讨论的连分数由 $y/y'$ 在 $x = 1$ 处的值给出. Oleinikov\index{Oleinikov, V. A.}\footnote{D.A.N. 166 (1966), 540--3.} 针对三阶线性微分方程获得了一些类似的定理；例如，他证明了若
\[F(x) = \sum_{n=0}^{\infty} \frac{(x/3)^{3n}}{n! [\lambda, n] [\mu, n]},\]
其中 $\lambda, \mu$ 为有理数，且 $\lambda + \mu, \lambda - 2\mu, \mu - 2\lambda$ 均非整数，则 $F(x), F'(x), F''(x)$ 满足\autoref{thm:ch11_1}的假设，从而对于每个非零代数数\index{algebraic number} $\alpha$，$F(\alpha), F'(\alpha), F''(\alpha)$ 代数无关. 而 Shidlovsky\index{Shidlovsky, A. B.}\footnote{Trudy Moskov. 18 (1968), 55--64.} 证明了一个引人注目的定理，大意是若
\[\Phi_k(x) = \sum_{n=0}^{\infty} x^{kn}/(n!)^k,\]
则对于任何非零代数数 $\alpha$ 以及任何 $r$，数\index{numbers} $\Phi_k^{(l)}(\alpha/k)$ （其中 $1 \leq l < k$，$1 \leq k \leq r$）代数无关. 显然，\autoref{thm:ch11_1} 也包含了 Lindemann\index{Lindemann, F.} 定理\index{Lindemann's theorem}.

\section{基本构造}\label{sec:ch11_2}

\autoref{thm:ch11_1} 的证明紧随前一章的论证，但确认 $\Delta(x)$ 不恒等于 0 已不再是一件简单的事. 该验证是 Shidlovsky\index{Shidlovsky, A. B.} 在该学科中的重大发现，将在下文的\autoref{lem:ch11_2}中给出.

我们用 $E_1(x), \dots, E_n(x)$ 表示如上所述在 $K(x)$ 上线性无关的 $E$-函数，并假设 $0 < \epsilon < 1$. 由 $\leqslant$ 或 $\gg$ 所隐含的常数将仅依赖于 $E$、 $f$ 中的系数以及 $\epsilon$. 我们用 $f(x)$ 表示一个在 $K$ 中系数不恒为 0 的多项式，使得对于 \eqref{eq:ch11_1} 中的所有 $f_{ij}$，$f f_{ij}$ 均为多项式.

\begin{mylemma}{}{ch11_1}
对于任何整数 $r \gg 1$，存在不全恒为 0 的多项式 $P_i(x)$ $(1 \leqslant i \leqslant n)$，其次数最多为 $r$，且在 $K$ 中的代数整数\index{algebraic integer}系数的大小最多为 $(r!)^{1+\epsilon}$，使得
\[\sum_{i=1}^{n} P_i(x) E_i(x) = \sum_{m=M}^{\infty} \rho_m x^m, \tag{2}\label{eq:ch11_2}\]
其中 $|\rho_m| < r! (m!)^{-1+\epsilon}$ 且
\[M = n(r+1) - 1 - [\epsilon r].\]
\end{mylemma}

\begin{proof}
设 $a_{ij}$ 为 $E_i(x)$ 中 $x^j/j!$ 的系数，并设 $b_{i0}, b_{i1}, \ldots$ 为如\autoref{sec:ch11_1}中与 $E_i$ 关联的整数序列. 由\autoref{ch:6}的\autoref{lem:ch6_1}，在 $K$ 中存在不全为 0 的代数整数\index{algebraic integer} $p_{ij}$ $(1 \le i \le n, 0 \le j \le r)$，使得
\[\sum_{i=1}^{n} \sum_{j=0}^{\min(r, m)} \binom{m}{j} a_{i, m-j} p_{ij} = 0 \quad (0 \le m < M), \tag{3}\label{eq:ch11_3}\]
并且实际上可以选定它们，使其大小最多为 $N^{\delta NM/(N-M)}$，其中 $N=n(r+1)$ 且 $\delta=(\epsilon/4n)^2$；因为在将 \eqref{eq:ch11_3} 乘以 $b_{1m}\dots b_{nm}$ 后，系数变为 $K$ 中的代数整数\index{algebraic integer}，其大小 $\leq 2^M M^{\frac{1}{2}\delta M}$，这在 $b$ 的定义性质中用 $\delta/(2n)$ 代替 $\epsilon$ 便很清楚. 我们如第10章引理 1 的证明中那样得出结论，多项式
\[P_i(x) = r! \sum_{j=0}^{r} p_{ij}(j!)^{-1} x^j \quad (1 \le i \le n)\]
具有所述的性质. 事实上，显然 \eqref{eq:ch11_2} 对于 $\rho_m = (r!/m!)\sigma_m$ 成立，其中 $\sigma_m$ 表示 \eqref{eq:ch11_3} 的左式，并且由于 $M < N < 2nr$ 且 $N - M > \epsilon r$，我们看到 $p_{ij}$ 的大小最多为 $r^{\frac{1}{2}\epsilon r}$，由此对于 $m \ge M$ 有 $|\sigma_m| < (m!)^{\epsilon}$，正如所求.
\end{proof}

\begin{mylemma}{}{ch11_2}
设 $P_{ij}(x)$ $(1 \le i \le n, j \ge 1)$ 递归定义为
\[P_{i1} = P_i, \quad P_{i,j+1} = f P'_{ij} + f \sum_{h=1}^{n} f_{hi} P_{hj}.\]
则以 $P_{ij}(x)$ 为第 $i$ 行第 $j$ 列元素的 $n$ 阶行列式 $\Delta(x)$ 不恒等于 0.
\end{mylemma}

\begin{proof}
反之，假设 $\Delta(x)$ 恒等于 0. 设 $k$ 为使得 $\Delta(x)$ 的前 $k$ 列在 $K(x)$ 上线性无关、但第 $k+1$ 列线性相关于这些列 of the integer. 我们用 $\mathbf{Q}$ 表示由 $\Delta(x)$ 的前 $k$ 列构成的矩阵，并用 $\mathbf{R}$ 和 $\mathbf{S}$ 分别表示由 $\mathbf{Q}$ 的前 $k$ 行和后 $n-k$ 行构成的矩阵. 我们假设（显然可以这样做）记号使得 $\mathbf{R}$ 是非奇异的，并着手证明 $\mathbf{S}\mathbf{R}^{-1}$ 的有理函数元素的分子和分母的次数均为 $\ll 1$，其中实际上隐含的常数仅依赖于 $f$. 这将足以确立引理；因为若用 $\mathbf{L}$ 表示第 $j$ 个元素为
\[L_{j} = \sum_{i=1}^{n} P_{ij} E_{i} \quad (1 \leq j \leq k),\]
的行向量，并令 $\mathbf{A} = (E_1, \dots, E_k)$，$\mathbf{B} = (E_{k+1}, \dots, E_n)$，我们有 $\mathbf{L} = \mathbf{A}\mathbf{R} + \mathbf{B}\mathbf{S}$，从而
\[\mathbf{L}\mathbf{R}^{-1} = \mathbf{A} + \mathbf{B}\mathbf{S}\mathbf{R}^{-1}. \tag{4}\label{eq:ch11_4}\]
但 $L_j$ 满足微分方程 $L_{j+1} = f L'_j$，因此 $\mathbf{L}$ 的每个元素在 $x=0$ 处均有至少 $M-n$ 阶的零点. 此外，$\mathbf{R}^{-1}$ 的每个元素均可表示为 $K(x)$ 中的有理函数，其分母为 $\det\mathbf{R}$，且由于后者是次数最多为 $kr+c$ （其中 $c \leqslant 1$）的多项式，因此 $\mathbf{L}\mathbf{R}^{-1}$ 的每个元素在 $x=0$ 处均有至少 $M-kr-n-c$ 阶的零点. 另一方面，鉴于所假设的 $E_1, \ldots, E_n$ 的线性无关性，\eqref{eq:ch11_4} 右侧的向量不能恒等于 0，且其元素在 $x=0$ 处的零点阶数（如果有的话）是有界的，该界独立于 $\mathbf{S}\mathbf{R}^{-1}$ 的元素的系数，因而特别地独立于 $r$. 现在 $k < n$，所以 $M-kr$ 随 $r$ 趋于无穷大；因此若 $r$ 足够大，我们便得到矛盾.

为了证明关于 $\mathbf{S}\mathbf{R}^{-1}$ 的断言，我们首先观察到存在一个以 $K(x)$ 为元素的 $k$ 阶方阵 $F$，使得对于 \eqref{eq:ch11_1} 的任何解 $\mathbf{y}$，向量 $\mathbf{Y} = \mathbf{y}\mathbf{Q}$ 满足微分方程 $\mathbf{Y}' = \mathbf{Y}F$. 事实上，若 $Y_j$ 表示 $\mathbf{Y}$ 的第 $j$ 个元素，则对于所有 $j < k$ 有 $Y_{j+1} = f Y'_j$，并且由 $k$ 的定义，$f Y'_k$ 可表示为 $Y_1, \ldots, Y_k$ 具有 $K(x)$ 中系数的线性组合. 现在设 $\mathbf{w}_1, \dots, \mathbf{w}_n$ 为在 $K$ 上线性无关的 \eqref{eq:ch11_1} 的幂级数解，并设 $\mathbf{W}$ 为以 $\mathbf{w}_j$ 为第 $j$ 行的 $n$ 阶方阵. 则 $\mathbf{W}\mathbf{Q}$ 的每一行都是 $\mathbf{Y}' = \mathbf{Y}F$ 的解；但这有最多 $k$ 个在 $K$ 上线性无关的解，因此存在一个具有 $K$ 中系数且秩为 $n-k$ 的 $n-k$ 乘 $n$ 矩阵 $\mathbf{M}$ 满足 $\mathbf{M}\mathbf{W}\mathbf{Q} = 0$. 用 $\mathbf{U}$ 和 $\mathbf{V}$ 分别表示由 $\mathbf{M}\mathbf{W}$ 的前 $k$ 列和后 $n-k$ 列构成的矩阵，我们有 $\mathbf{U}\mathbf{R} + \mathbf{V}\mathbf{S} = 0$. 因为 $\mathbf{R}$ 是非奇异的且 $\mathbf{M}\mathbf{W}$ 的秩为 $n-k$，由此可知 $\mathbf{V}$ 是非奇异的，从而 $\mathbf{S}\mathbf{R}^{-1} = -\mathbf{V}^{-1}\mathbf{U}$. 显然 $\mathbf{V}^{-1}\mathbf{U}$ 的元素是 $\mathbf{W}$ 的元素的具有 $K$ 中系数的有理函数，且其分子和分母的次数是有界的，该界独立于 $r$. 因此它们可以表示为 $\mathbf{W}$ 的元素的某些在 $K(x)$ 上线性无关 of certain monomials 的线性型的商，线性型中的系数为 $K(x)$ 中的有理函数，这些有理函数的分子和分母的次数再次有界，其界独立于 $r$. 由于 $\mathbf{S}\mathbf{R}^{-1}$ 从而 $\mathbf{V}^{-1}\mathbf{U}$ 的元素实际上在 $K(x)$ 中，它们必须由这些系数的商给出，由此断言成立.
\end{proof}

\section{进一步的引理}\label{sec:ch11_3}

我们现在建立第10章引理 3 和引理 4 的类似命题. 此处的论证将紧随先前的对应论证，因此我们将较为简短.

We shall signify by $\alpha$ an element of $K$ with $\alpha f(\alpha) \neq 0$. We denote by $c_1, c_2, \ldots$ positive numbers\index{numbers} which may depend on $\alpha, \epsilon$ and the coefficients in the $E$'s and $f$'s only.

\begin{mylemma}{}{ch11_3}
存在不超过 $\epsilon r + c_1$ 的相异下标 $J(j)$ $(1 \le j \le n)$，使得以 $P_{i, J(j)}(x)$ 为第 $i$ 行第 $j$ 列元素的行列式在 $x = \alpha$ 处不为 0.
\end{mylemma}

\begin{proof}
我们首先注意到，若将第一行替换为 $E_1^{-1}L_j$（其中 $j=1,2,\ldots,n$ 且 $L_j$ 定义见\autoref{lem:ch11_2}的证明），$\Delta(x)$ 保持不变. 因此 $\Delta(x)$ 在 $x=0$ 处有至少 $M-c_2$ 阶的零点，并且由于它是次数最多为 $nr+c_3$ 的多项式，因此存在一个不超过
\[nr + c_3 - (M - c_2) \leqslant \epsilon r + c_4\]
的非负整数 $t$ 使得 $\Delta^{(t)}(\alpha) \neq 0$；我们假设 $t$ 是极小化选取的.

我们现在通过第10章的 (3) 引入 $w_1, \ldots, w_n$ 的线性型. 通过对该章的 (4) 应用算子 $f d/dx$ 共 $t$ 次，并将每次出现的 $w'_i$ 替换为 \eqref{eq:ch11_1} 右侧在 $y_j = w_j$ 处的值，我们得到
\[w_i(f(\alpha))^t \Delta^{(t)}(\alpha) = \sum_{j=1}^{n+t} W_j F_{ij}(\alpha) \quad (1 \leq i \leq n),\]
其中 $F_{ij}$ 表示 $x$ 的多项式，由 $f$、$\Delta$ 及其导数的线性组合给出. 因此，当 $x = \alpha$ 时，线性型
\[W_j \quad (1 \leqslant j \leqslant n+t)\]
包含一组 $n$ 个线性无关的型，并且引理对于由相关联的下标给出的 $J(j)$ 成立.
\end{proof}

\begin{mylemma}{}{ch11_4}
在 $K$ 中存在大小最多为 $(r!)^{1+18\epsilon}$ 的代数整数\index{algebraic integer} $q_{ij}$ $(1 \le i, j \le n)$，它们构成非零行列式并满足
\[\left|\sum_{i=1}^{n} q_{ij} E_i(\alpha)\right| < (r!)^{-n+1+16\epsilon n} \quad (1 \leq j \leq n). \tag{5}\label{eq:ch11_5}\]
\end{mylemma}

\begin{proof}
设 $l$ 为正整数，使得 $l\alpha$ 以及 $l f$ 和所有 $l f f_{ij}$ 的系数均为代数整数\index{algebraic integer}. 我们着手证明数\index{numbers}
\[q_{ij} = l^{r+(m+1)J(j)} P_{i, J(j)}(\alpha) \quad (1 \leq i, j \leq n),\]
其中 $m$ 表示 $f f_{ij}$ 和 $f$ 的次数的最大值，具有所求的性质. 首先，显然 $l^j P_{ij}$ 具有代数整数\index{algebraic integer}系数且次数最多为 $r+mj$. 因此 $q_{ij}$ 是代数整数\index{algebraic integer}，并且由\autoref{lem:ch11_3}，它们构成非零行列式. 此外，通过归纳法易证 $l^j P_{ij}$ 系数的大小最多为 $(r+mj)^{2j} c_5^j (r!)^{1+\epsilon}$，且由于 $J(j)$ 不超过 $\epsilon r+c_1$，这给出了 $q_{ij}$ 的大小的所需估计.

余下需证明 \eqref{eq:ch11_5}. 用 $\Phi(x)$ 表示 \eqref{eq:ch11_2} 的左式，显然 \eqref{eq:ch11_5} 左侧的和除了因子 $l^{r+(m+1)J}$ 外，由在 $x=\alpha$ 处求值的 $(f d/dx)^{J-1}\Phi$ 给出，其中 $J=J(j)$. 但是，再次通过归纳法，我们看到这是 $\Phi^{(h)}(\alpha)$（其中 $h=0,1,\dots,J-1$）的线性型，其系数的绝对值最多为 $(c_6 J)^{2J}$. 因此只需证明
\[|\Phi^{(h)}(\alpha)| < (r!)^{-n+1+8\epsilon n} \quad (0 \le h < J).\]
现在由\autoref{lem:ch11_1}，我们得到
\[|\Phi^{(h)}(\alpha)| < r! \sum_{m=M}^{\infty} (m!)^{-1+\epsilon} ((m-h)!)^{-1} |\alpha|^{m-h},\]
且右侧的和最多为
\[h! \sum_{m=M}^{\infty} (m!)^{-1+\epsilon} 2^m |\alpha|^{m-h} \leqslant h! c_7^M (M!)^{-1+\epsilon}.\]
由于 $h < \epsilon r + c_1$ 且 $M \leq 2nr$，我们有 $h! \leq (r!)^{3\epsilon}$ 并且
\[M! \geqslant (2nr)^{-\epsilon r} (r!)^n \geqslant (r!)^{n-3\epsilon}.\]
所需的估计立即成立.
\end{proof}

\section{主定理的证明}\label{sec:ch11_4}

假设 $E_1(\alpha), \ldots, E_n(\alpha)$ 代数相关. 则它们满足方程 $P(E_1, \ldots, E_n) = 0$，其中 $P$ 是一个代数系数不全为 0 的多项式. 我们用 $c$ 表示 $P$ 的次数，并且我们在不失一般性的情况下，可以假设 $P$ 的系数为 $K$ 中的代数整数\index{algebraic integer}. $K$ 的次数将用 $d$ 表示，并且我们假设 $0 < \epsilon < 1$. 此外，我们用 $m$ 表示一个整数，使得二项式系数
\[k = \binom{m+n+c}{n}, \quad l = \binom{m+n}{n}\]
满足 $k - l < l/(2d)$；由于当 $m \to \infty$ 时 $k$ and $l$ are asymptotic to $m^n/n!$ as $m \to \infty$（通过将它们表示为 $m$ 的 $n$ 次多项式便很容易看出），上述不等式对于所有足够大的 $m$ 确实成立. 在下文中，由 $\ll$ 或 $\gg$ 隐含的常数将仅依赖于 $\alpha, \epsilon, m$ 以及 $E$、 $f$ 和 $P$ 中的系数.

现在设 $\mathscr{E}_1, \ldots, \mathscr{E}_k$ 为 $E$-函数 $E_1^{j_1} \dots E_n^{j_n}$，其中 $j_1, \ldots, j_n$ run through all non-negative integers with $j_1 + \ldots + j_n < m + c$. 那么显然，$\mathscr{E}_1, \ldots, \mathscr{E}_k$ satisfy a further system of linear differential equations of the form \eqref{eq:ch11_1}, where the new coefficients are given by linear combinations of the $f$'s;此外，鉴于关于 $E_1(x), \ldots, E_n(x)$ 代数无关性的假设，$\mathscr{E}_1, \ldots, \mathscr{E}_k$ 在 $K(x)$ 上线性无关. 我们从\autoref{sec:ch11_2}和\autoref{sec:ch11_3}得出结论，对于任何整数 $r \gg 1$，在 $K$ 中存在具有\autoref{lem:ch11_4}中所引性质的代数整数\index{algebraic integer} $q_{ij}$ ($1 \le i, j \le k$)，其中以 $\mathscr{E}_1, \dots, \mathscr{E}_k$ 代替 $E_1, \dots, E_n$. 对于每个 set of non-negative integers $j_1, \dots, j_n$ with $j_1 + \ldots + j_n < m$ we write
\[E_1^{j_1} \dots E_n^{j_n} P(E_1, \dots, E_n) = \sum_{i=1}^k p_{ij} \mathscr{E}_i,\]
其中 $p_{ij}$ 为 $P$ 的系数或 0，且 $j = j(j_1, \ldots, j_n)$ takes the values $1, 2, \ldots, l$. 于是在右侧我们有 $l$ 个在 $K$ 上线性无关的 $\mathscr{E}_1, \dots, \mathscr{E}_k$ 的线性型，它们在 $x = \alpha$ 处均消失. 由于 $q_{ij}$ 的行列式不为 0，因此可知存在 $k - l$ of the forms
\[\Phi_j = \sum_{i=1}^k q_{ij} \mathscr{E}_i \quad (1 \leqslant j \leqslant k)\]
which together with the latter make up a linearly independent set;在不失一般性的情况下，我们可以假设它们由 $\Phi_{l+1}, \ldots, \Phi_k$ 给出. 我们还可以假设（显然可以这样做） $\mathscr{E}_1(\alpha) \neq 0$.

我们现在比较当 $j \le l$ 时以 $p_{ij}$ 为第 $i$ 行第 $j$ 列元素、当 $j > l$ 时以 $q_{ij}$ 在该位置的 $k$ 阶行列式 $D$ 的估计. 显然 $D$ 是 $K$ 中的非零代数整数\index{algebraic integer}，并且由于 $p_{ij} \ll 1$，它的大小为 $\ll (r!)^{(1+18\epsilon)(k-l)}$；因此
\[|D| \gg (r!)^{-(1+18\epsilon)(k-l)d} \geqslant (r!)^{-(1+18\epsilon)l/2}.\]
另一方面，如果将第一行在 $j \le l$ 时替换为 0、在 $j > l$ 时替换为 $\mathscr{E}_1^{-1}(\alpha)\Phi_j$，则 $D$ 保持不变. 此外，由\autoref{lem:ch11_4}，后者的元素为 $\ll (r!)^{-k+1+16\epsilon k}$；因此
\[|D| \ll (r!)^{(1+18\epsilon)(k-l-1)-k+1+18\epsilon k} \leqslant (r!)^{-l+32\epsilon k}.\]
但是 $k < \frac{3}{2} l$，因此若 $\epsilon < 1/(100k)$ 并且 $r$ 足够大，我们便得到矛盾. 这证明了定理.

在 Shidlovsky\index{Shidlovsky, A. B.} 做出根本性发现之后，该领域的研究主要集中于为特定类别的 $E$-函数确立\autoref{thm:ch11_1}及其推广的函数论假设，正如第 1 节中所指出的，这在许多引人注目的情形中确实已经实现. 人们还针对获取如上所述 $P(E_1, \ldots, E_n)$ 型式的正下界开展了研究，实际上已经建立了一个形如 $C h^{-c}$ 的估计，其中 $h$ 表示 $P$ 的系数大小的最大值，而 $C, c$ 为不依赖于 $h$ 的正数\index{numbers}；但这里的 $c$ 随 $n$ 的增加而迅速增大\footnote{参见 Lang\index{Lang, S.}（文献目录，第一部著作）. } . 该学科中主要的悬而未决的问题是将该理论推广到比 $E$-函数更广泛的解析函数类，此处的任何进展都将极具意义.
